\documentclass[10pt, a4paper]{article}

% ============================================================================
% PAQUETES ESTÁNDAR 100% COMPATIBLES CON OVERLEAF (pdfLaTeX) Y TEXLIVE LOCAL
% ============================================================================
\usepackage[utf8]{inputenc}
\usepackage[T1]{fontenc}
\usepackage[top=1.8cm, bottom=2.0cm, left=2.0cm, right=2.0cm, headheight=14pt]{geometry}

\usepackage{amsmath, amssymb, amsfonts}
\usepackage{booktabs}
\usepackage{tabularx}
\usepackage{array}
\usepackage{enumitem}
\usepackage{xcolor}
\usepackage{fancyhdr}
\usepackage{titlesec}
\PassOptionsToPackage{hyphens}{url}
\usepackage[colorlinks=true, linkcolor=blue!70!black, urlcolor=blue!70!black, citecolor=blue!70!black]{hyperref}

% ============================================================================
% VARIABLES EDITABLES ANTES DE LA ENTREGA FINAL (DOMINGO 18 DE OCTUBRE)
% Solo cambia el contenido entre llaves {} de estas dos líneas en Overleaf:
% ============================================================================
\newcommand{\CommitSHA}{\texttt{[PENDIENTE: Pegar SHA del commit final al congelar código (EID-14)]}}
\newcommand{\DuracionVideo}{\texttt{[PENDIENTE: Indicar minutos:segundos exactos (EID-20)]} (Rango: 4:00 a 6:00 min)}

% ============================================================================
% ESTILO Y FORMATO ACADÉMICO
% ============================================================================
\definecolor{uctblue}{RGB}{0, 56, 118}
\definecolor{lightgray}{RGB}{245, 247, 250}

\pagestyle{fancy}
\fancyhf{}
\fancyhead[L]{\small\textbf{MAT1188 -- Cálculo Intermedio} $\cdot$ Ficha de Entrega EID}
\fancyhead[R]{\small\textbf{Grupo 3 (G03)} $\cdot$ Semestre 2026-2}
\fancyfoot[C]{\small Página \thepage}
\renewcommand{\headrulewidth}{0.4pt}
\renewcommand{\footrulewidth}{0.4pt}

\titleformat{\section}
  {\color{uctblue}\normalfont\large\bfseries}
  {\thesection.}{0.5em}{}[\vspace{-0.3em}\textcolor{uctblue}{\titlerule[0.6pt]}]
\titlespacing*{\section}{0pt}{1.0em}{0.5em}

\titleformat{\subsection}
  {\color{uctblue!85!black}\normalfont\normalsize\bfseries}
  {\thesubsection.}{0.5em}{}
\titlespacing*{\subsection}{0pt}{0.7em}{0.35em}

\setlength{\parindent}{0pt}
\setlength{\parskip}{0.4em}
\renewcommand{\arraystretch}{1.22}
\newcolumntype{Y}{>{\raggedright\arraybackslash}X}

\begin{document}

% ============================================================================
% ENCABEZADO OFICIAL (FORMATO PAUTA EID MAT1188)
% ============================================================================
\begin{center}
    {\LARGE\bfseries\color{uctblue} Ficha de Entrega -- Evaluación Integrada de Desempeño (EID)}\\[0.25em]
    {\large\bfseries Introducción a la Integral con Manim Community}\\[0.15em]
    {\normalsize Universidad Católica de Temuco $\cdot$ Ingeniería Civil Informática (2.\textsuperscript{o} Año)}
\end{center}

\vspace{0.2em}
\noindent
\fcolorbox{uctblue}{lightgray}{%
\begin{minipage}{\dimexpr\textwidth-2\fboxsep-2\fboxrule\relax}
\vspace{0.25em}
\textbf{Curso:} MAT1188 -- Cálculo Intermedio (Semestre 2026-2) \hfill \textbf{Fecha:} 18 de octubre de 2026\\[0.15em]
\textbf{Grupo:} Grupo 3 (\texttt{G03})\\[0.15em]
\textbf{Tema asignado:} Eje 3 $\cdot$ Aproximación y refinamiento -- \textit{«De pocos rectángulos a una mejor aproximación»}\\[0.15em]
\textbf{Pregunta central:} \textit{«¿Por qué una suma de Riemann puede aproximar cada vez mejor una región cuando aumentamos el número de rectángulos?»}\\[0.25em]
\textbf{Integrantes del equipo:}
\begin{itemize}[leftmargin=1.5em, itemsep=1pt, topsep=2pt]
    \item \textbf{Denys Sofía Rodríguez Noguera} (\href{https://github.com/sofiarodrigueznog}{\texttt{@sofiarodrigueznog}}) -- \textit{Líder del equipo / Revisión e integración del video}
    \item \textbf{Paulo Villalobos Aravena} (\href{https://github.com/notKechai}{\texttt{@notKechai}}) -- \textit{Matemática y explicación teórica}
    \item \textbf{Antonio Ignacio Lara Fuentealba} (\href{https://github.com/alara2024uct}{\texttt{@alara2024uct}}) -- \textit{Visualización y programación en Manim}
    \item \textbf{Renato Misrrain Carrasco Oviedo} (\href{https://github.com/RCarrascoO}{\texttt{@RCarrascoO}}) -- \textit{Apoyo matemático, visualización y auditoría}
\end{itemize}
\vspace{0.1em}
\end{minipage}%
}

% ============================================================================
% 1. IDENTIFICACIÓN Y METADATOS DE ENTREGA
% ============================================================================
\section{Identificación Técnica y Localización Documental (Criterio V8)}

\begin{tabularx}{\textwidth}{@{} >{\bfseries}p{5.2cm} Y @{}}
\toprule
Campo Requerido & Información Oficial de Entrega \\
\midrule
Código y Asignatura & MAT1188 -- Cálculo Intermedio \\
Número de Grupo & Grupo 3 (\texttt{G03}) \\
Integrantes & Denys Sofía Rodríguez Noguera, Paulo Villalobos Aravena, Antonio Ignacio Lara Fuentealba y Renato Misrrain Carrasco Oviedo \\
Repositorio Público de GitHub & \url{https://github.com/sofiarodrigueznog/EID-Grupo-3} \\
SHA del Commit Entregado & \CommitSHA \\
Archivo de Video Entregado & \texttt{MAT1188\_EID\_G03\_2026.mp4} \\
Duración Exacta del Video & \DuracionVideo \\
Archivo de Ficha PDF & \texttt{MAT1188\_EID\_G03\_2026.pdf} \\
Código de Ética & Entregado en Blackboard (\texttt{MAT1188\_ETICA\_G03\_2026.pdf}) el 06/10/2026 \\
\bottomrule
\end{tabularx}

\subsection*{Mapa de Localización de Evidencias en el Repositorio (V8d, V8e)}
Todos los archivos que respaldan el video y la defensa oral se localizan directamente desde la raíz del repositorio:

\begin{tabularx}{\textwidth}{@{} p{4.5cm} p{4.1cm} Y @{}}
\toprule
\textbf{Documento / Recurso} & \textbf{Ruta en el Repositorio} & \textbf{Contenido y Propósito} \\
\midrule
\textbf{Guion y Notas Matemáticas} & \href{https://github.com/sofiarodrigueznog/EID-Grupo-3/blob/main/docs/guion.md}{\texttt{docs/guion.md}} & Guion sincronizado por escena, tiempos, hablantes y cálculos que respaldan la locución. \\
\textbf{Marco Teórico} & \href{https://github.com/sofiarodrigueznog/EID-Grupo-3/blob/main/docs/teoria.md}{\texttt{docs/teoria.md}} & Definiciones de partición, puntos muestra, suma de Riemann, refinamiento e integral definida. \\
\textbf{Ejemplo Desarrollado} & \href{https://github.com/sofiarodrigueznog/EID-Grupo-3/blob/main/docs/ejemplo.md}{\texttt{docs/ejemplo.md}} & Desarrollo paso a paso para $n=4, 8, 16, 32$, fórmula cerrada $S_n$ y análisis del error. \\
\textbf{Fuentes y Referencias} & \href{https://github.com/sofiarodrigueznog/EID-Grupo-3/blob/main/docs/referencias.md}{\texttt{docs/referencias.md}} & Bibliografía completa localizable (apunte base, libros con ISBN y documentación de Manim). \\
\textbf{Declaración de Uso de IA} & \href{https://github.com/sofiarodrigueznog/EID-Grupo-3/blob/main/docs/declaracion_ia.md}{\texttt{docs/declaracion\_ia.md}} & Desglose por integrante de herramientas de IA utilizadas, propósito y verificación humana. \\
\textbf{Parámetros Centralizados} & \href{https://github.com/sofiarodrigueznog/EID-Grupo-3/blob/main/scenes/config.py}{\texttt{scenes/config.py}} & Configuración única ($f$, $[a,b]$, $n$, punto medio) importada por las tres escenas. \\
\textbf{Auditoría y Ejecución} & \href{https://github.com/sofiarodrigueznog/EID-Grupo-3/blob/main/verify.py}{\texttt{verify.py}} / \href{https://github.com/sofiarodrigueznog/EID-Grupo-3/blob/main/render.py}{\texttt{render.py}} & Verificación racional exacta (\texttt{Fraction}), suite de tests (\texttt{tests/}) y runner multiplataforma. \\
\bottomrule
\end{tabularx}

% ============================================================================
% 2. NÓMINA DE INTEGRANTES Y RESPONSABILIDADES
% ============================================================================
\section{Nómina de Integrantes, Roles y Aportes Efectivos}

\begin{tabularx}{\textwidth}{@{} p{4.0cm} p{3.5cm} p{3.1cm} Y @{}}
\toprule
\textbf{Integrante} & \textbf{Cuenta GitHub} & \textbf{Rol Principal} & \textbf{Aporte Efectivo al Proyecto} \\
\midrule
\textbf{Denys Sofía Rodríguez Noguera} \textit{(Líder)} & \href{https://github.com/sofiarodrigueznog}{\texttt{@sofiarodrigueznog}} & Revisión e integración del video & Coordinación general, arquitectura del repositorio, workflows de CI, montaje del MP4 final, locución y entrega oficial. \\
\textbf{Paulo Villalobos Aravena} & \href{https://github.com/notKechai}{\texttt{@notKechai}} & Matemática y explicación & Redacción de \texttt{teoria.md} y \texttt{ejemplo.md}, deducción analítica de la fórmula cerrada $S_n$, guion y locución. \\
\textbf{Antonio Ignacio Lara Fuentealba} & \href{https://github.com/alara2024uct}{\texttt{@alara2024uct}} & Visualización y código & Parametrización en \texttt{scenes/config.py}, tests unitarios (\texttt{test\_config.py}), programación de Escenas 1 y 2, y locución. \\
\textbf{Renato Misrrain Carrasco Oviedo} & \href{https://github.com/RCarrascoO}{\texttt{@RCarrascoO}} & Apoyo matemática y visualización & Script de auditoría racional (\texttt{verify.py}), runner (\texttt{render.py}), programación de Escena 3, locución y ficha técnica. \\
\bottomrule
\end{tabularx}

% ============================================================================
% 3. FUNDAMENTOS MATEMÁTICOS Y DESARROLLO DEL EJEMPLO
% ============================================================================
\section{Fundamentos Matemáticos y Desarrollo del Ejemplo (Criterios V1, V2, V3)}

\subsection{Definición Formal y Glosario de Símbolos (V1a, V1b)}
Sea $f:[a,b] \to \mathbb{R}$ una función definida en el intervalo cerrado $[a,b]$. Una \textbf{partición regular} divide $[a,b]$ en $n \in \mathbb{N}$ subintervalos $[x_{i-1}, x_i]$ de igual ancho mediante los nodos $x_i = a + i\Delta x$ ($i=0,1,\dots,n$). La \textbf{suma de Riemann} se define como:
\[
S_n = \sum_{i=1}^{n} f(x_i^*)\,\Delta x
\]
\begin{itemize}[leftmargin=1.5em, itemsep=1pt, topsep=2pt]
    \item $[a,b] = [0,2]$: Intervalo cerrado y acotado de integración; $n$: cantidad de subintervalos (y de rectángulos).
    \item $\Delta x = \dfrac{b-a}{n} = \dfrac{2}{n}$: Ancho uniforme de cada subintervalo (base de cada rectángulo y norma de la partición $\|P\|$).
    \item $x_i^* \in [x_{i-1}, x_i]$: Punto muestra del subintervalo $i$. Empleamos el \textbf{punto medio}: $x_i^* = \dfrac{x_{i-1}+x_i}{2} = a + \left(i - \dfrac{1}{2}\right)\Delta x$.
    \item $f(x_i^*)$: Altura del $i$-ésimo rectángulo; $A_i = f(x_i^*)\,\Delta x$: Área del $i$-ésimo rectángulo; $\sum_{i=1}^n$: suma finita de las $n$ áreas.
\end{itemize}

\subsection{Condiciones Necesarias y Distinción de Conceptos (V1c, V1d)}
\begin{enumerate}[leftmargin=1.5em, itemsep=2pt, topsep=2pt]
    \item \textbf{Continuidad e integrabilidad:} Como $f(x) = x^2 + 1$ es un polinomio, es \textbf{continua} en todo $[0,2]$. La continuidad en un intervalo cerrado garantiza que $f$ es integrable según Riemann y que el límite $\lim_{n\to\infty} S_n$ existe y es único, sin importar la elección de $x_i^*$.
    \item \textbf{Condición de no-negatividad ($f(x) \geq 0$):} En general, la integral definida representa un \textit{área neta con signo}. Como en nuestro ejemplo $f(x) = x^2 + 1 \geq 1 > 0$ en todo $[0,2]$, las alturas $f(x_i^*)$ son estrictamente positivas y tanto $S_n$ como la integral representan directamente el \textbf{área geométrica} bajo la curva.
    \item \textbf{Distinción conceptual:} \textbf{(i)} $S_n$ es una aproximación discreta para un $n$ finito; \textbf{(ii)} el refinamiento ($n \to \infty \iff \Delta x \to 0$) es el proceso de paso al límite; \textbf{(iii)} la integral definida $\int_a^b f(x)\,dx = \lim_{n\to\infty} \sum_{i=1}^n f(x_i^*)\,\Delta x$ es el valor real exacto límite.
\end{enumerate}

\subsection{Desarrollo Analítico del Ejemplo y Deducción de la Fórmula Cerrada (V2a--d)}
\begin{itemize}[leftmargin=1.5em, itemsep=3pt, topsep=2pt]
    \item \textbf{Valor exacto por el Teorema Fundamental del Cálculo:}
    \[
    \int_0^2 (x^2 + 1)\,dx = \left[ \frac{x^3}{3} + x \right]_0^2 = \frac{2^3}{3} + 2 = \frac{8}{3} + 2 = \frac{14}{3} \approx 4.666667
    \]
    \item \textbf{Cálculo explícito para $n=4$:} Con $\Delta x = \frac{2}{4} = \frac{1}{2}$, los puntos medios son $x_1^* = \frac{1}{4}, x_2^* = \frac{3}{4}, x_3^* = \frac{5}{4}, x_4^* = \frac{7}{4}$. Evaluando:
    \[
    S_4 = \frac{1}{2}\left[\left(\frac{1}{16}+1\right) + \left(\frac{9}{16}+1\right) + \left(\frac{25}{16}+1\right) + \left(\frac{49}{16}+1\right)\right] = \frac{1}{2}\left[\frac{84}{16} + 4\right] = \frac{1}{2}\cdot\frac{37}{4} = \frac{37}{8} = 4.625
    \]
    \item \textbf{Deducción algebraica de la fórmula cerrada general $S_n$:}
    Para cualquier $n \in \mathbb{N}$, con $\Delta x = \frac{2}{n}$ y $x_i^* = \left(i - \frac{1}{2}\right)\frac{2}{n} = \frac{2i-1}{n}$:
    \[
    S_n = \sum_{i=1}^{n} \left[\left(\frac{2i-1}{n}\right)^2 + 1\right]\frac{2}{n} = \frac{2}{n}\left[\frac{1}{n^2}\sum_{i=1}^{n}(4i^2 - 4i + 1) + \sum_{i=1}^{n} 1\right]
    \]
    Aplicando las fórmulas de sumas de potencias $\sum_{i=1}^n i^2 = \frac{n(n+1)(2n+1)}{6}$, $\sum_{i=1}^n i = \frac{n(n+1)}{2}$ y $\sum_{i=1}^n 1 = n$, la suma de cuadrados de impares se reduce a $\sum_{i=1}^n (2i-1)^2 = \frac{n(4n^2-1)}{3}$. Sustituyendo:
    \[
    S_n = \frac{2}{n}\left[\frac{1}{n^2}\cdot\frac{n(4n^2-1)}{3} + n\right] = \frac{2(4n^2-1)}{3n^2} + 2 = \frac{8n^2 - 2 + 6n^2}{3n^2} = \frac{14n^2 - 2}{3n^2} = \frac{14}{3} - \frac{2}{3n^2}
    \]
\end{itemize}

\subsection{Tabla Comparativa de Refinamiento y Análisis del Error (V2d, V3a)}

\begin{center}
\setlength{\tabcolsep}{8pt}
\begin{tabular}{@{} c c c c c c @{}}
\toprule
\textbf{Subintervalos $n$} & \textbf{Ancho $\Delta x$} & \textbf{$S_n$ exacta} & \textbf{$S_n$ decimal} & \textbf{Error exacto $\left|\frac{14}{3} - S_n\right|$} & \textbf{Error decimal} \\
\midrule
$\mathbf{4}$ & $1/2 = 0.5000$ & $37/8$ & $4.625000$ & $1/24$ & $0.041667$ \\
$\mathbf{8}$ & $1/4 = 0.2500$ & $149/32$ & $4.656250$ & $1/96$ & $0.010417$ \\
$\mathbf{16}$ & $1/8 = 0.1250$ & $597/128$ & $4.664063$ & $1/384$ & $0.002604$ \\
$\mathbf{32}$ & $1/16 = 0.0625$ & $2389/512$ & $4.666016$ & $1/1536$ & $0.000651$ \\
\midrule
$\mathbf{n \to \infty}$ & $\Delta x \to 0$ & $\mathbf{14/3}$ & $\mathbf{4.666667}$ & $\mathbf{0}$ & $\mathbf{0.000000}$ \\
\bottomrule
\end{tabular}
\end{center}

\begin{itemize}[leftmargin=1.5em, itemsep=2pt, topsep=2pt]
    \item \textbf{Justificación geométrica del signo del error (Convexidad):} Dado que $f''(x) = 2 > 0$, la parábola es estrictamente convexa (cóncava hacia arriba). La recta tangente en el punto medio queda por debajo de la curva, por lo que en cada subintervalo el área que falta supera levemente al área que sobra ($S_n < \frac{14}{3}$).
    \item \textbf{Orden de convergencia cuadrático $O(1/n^2)$:} El error absoluto es $E_n = \frac{2}{3n^2}$. Al duplicar $n$, el error se divide exactamente por $4$ ($\frac{1}{24} \to \frac{1}{96} \to \frac{1}{384} \to \frac{1}{1536}$), logrando una convergencia más rápida que los extremos izquierdo o derecho ($O(1/n)$).
\end{itemize}

\subsection{Respuesta a la Pregunta Central (V3c, V3d, V3e)}
\fcolorbox{uctblue!60}{lightgray}{%
\begin{minipage}{\dimexpr\textwidth-2\fboxsep-2\fboxrule\relax}
Una suma de Riemann aproxima cada vez mejor la región porque, al mantener fijo el intervalo $[a,b]$ y aumentar la cantidad de rectángulos $n$, el ancho $\Delta x = \frac{b-a}{n}$ disminuye hacia cero. En tramos tan angostos, una función continua varía muy poco dentro de cada subintervalo, por lo que la altura constante $f(x_i^*)$ evaluada en el punto medio representa con mayor fidelidad el comportamiento local de la curva, reduciendo el área sobrante y faltante. Analíticamente, a partir de $S_n = \frac{14}{3} - \frac{2}{3n^2}$, al tomar el límite cuando $n \to \infty$ la discrepancia $\frac{2}{3n^2} \to 0$, demostrando que las sumas de Riemann convergen al valor exacto de la integral definida $\int_0^2 (x^2+1)\,dx = \frac{14}{3}$.
\end{minipage}%
}

% ============================================================================
% 4. ESCENAS DE MANIM Y ESTRUCTURA DE PARTICIPACIÓN EN VIDEO
% ============================================================================
\section{Escenas de Manim y Participación en el Video (Criterios V4, V5, V6)}

El proyecto emplea \textbf{Manim Community v0.22.0} sobre \textbf{Python 3.12}. Todas las escenas importan sus datos desde \texttt{scenes/config.py}. Los cuatro integrantes participan con su propia voz explicando contenido matemático sustantivo y son identificados en pantalla:

\begin{tabularx}{\textwidth}{@{} p{4.6cm} p{2.1cm} p{2.7cm} Y @{}}
\toprule
\textbf{Escena / Archivo} & \textbf{Tramo Base} & \textbf{Locutor(a)} & \textbf{Contenido Matemático y Visualización} \\
\midrule
\textbf{Presentación} & Inicio & Sofía Rodríguez & Identificación del curso, Grupo 3, integrantes y tema asignado. \\
\midrule
\textbf{Escena 1: Construcción}\newline{\small\texttt{escena1\_construccion.py}}\newline({\small\texttt{Escena1Construccion}}) & 0:00 -- 0:09\newline 0:09 -- 0:22\newline 0:22 -- 0:29.5 & Paulo Villalobos\newline Sofía Rodríguez\newline Renato Carrasco & \textbf{Paulo:} Gráfica de $f(x)=x^2+1 \geq 0$ en $[0,2]$ y Pregunta Central.\newline \textbf{Sofía:} Partición $n=4$, $\Delta x = 0.5$ y punto medio $x_i^* = \frac{x_{i-1}+x_i}{2}$.\newline \textbf{Renato:} Construcción de los 4 rectángulos, cálculo de $A_1 = 0.53125$ y $S_4 = \frac{37}{8} = 4.625$. \\
\midrule
\textbf{Escena 2: Refinamiento}\newline{\small\texttt{escena2\_refinamiento.py}}\newline({\small\texttt{Escena2Refinamiento}}) & 0:00 -- 0:17\newline 0:17 -- 0:35.2 & Renato Carrasco\newline Sofía Rodríguez & \textbf{Renato:} Transición $n=4 \to 8$, reducción de $\Delta x$ a $0.25$ y mejora a $S_8 = 4.656250$.\newline \textbf{Sofía:} Paso a $n=16$ y $n=32$ ($S_{32} = 4.666016$), tabla comparativa, cálculo de $\int_0^2(x^2+1)dx = \frac{14}{3}$ y reducción del error a $0.000651$. \\
\midrule
\textbf{Escena 3: Límite e Integral}\newline{\small\texttt{escena3\_integral.py}}\newline({\small\texttt{Escena3Integral}}) & 0:00 -- 0:18\newline 0:18 -- 0:23\newline 0:23 -- 0:27.4 & Antonio Lara\newline Paulo Villalobos\newline Todos (Cierre) & \textbf{Antonio:} Proceso $n\to\infty \implies \Delta x\to 0$, paso formal $\lim_{n\to\infty}\sum f(x_i^*)\Delta x = \int_a^b f(x)dx$ y distinción entre representación finita ($n=32$) y límite.\newline \textbf{Paulo:} Análisis del error analítico $\frac{2}{3n^2} \to 0$.\newline \textbf{Todos:} Respuesta síntesis a la Pregunta Central y créditos. \\
\bottomrule
\end{tabularx}

% ============================================================================
% 5. ARQUITECTURA DE PARÁMETROS Y VARIACIONES (DEFENSA ORAL)
% ============================================================================
\section{Trazabilidad del Código y Respuesta ante Variaciones (Criterios D4 y D5)}

En \texttt{scenes/config.py} se centralizan: \texttt{f(x)}, \texttt{a = 0.0}, \texttt{b = 2.0}, \texttt{n\_values = [4, 8, 16, 32]} y \texttt{sample\_type = "center"} (con soporte para \texttt{"left": 0.0}, \texttt{"center": 0.5}, \texttt{"right": 1.0} en \texttt{sample\_x(i, n)}).
\begin{itemize}[leftmargin=1.5em, itemsep=2pt, topsep=2pt]
    \item \textbf{Variación de punto muestra (\texttt{sample\_type = "left"} o \texttt{"right"}):} Al ser $f(x)=x^2+1$ creciente en $[0,2]$, el extremo izquierdo subestima ($S_n < \frac{14}{3}$) y el extremo derecho sobreestima ($S_n > \frac{14}{3}$), con convergencia lineal $O(1/n)$ hacia el mismo límite $\frac{14}{3}$.
    \item \textbf{Variación del intervalo (p. ej. \texttt{b = 4.0}):} Para un $n$ fijo, $\Delta x = \frac{4}{n}$ se duplica y el error crece, requiriendo mayor $n$ para igual precisión.
\end{itemize}

% ============================================================================
% 6. FUENTES Y REFERENCIAS CONSULTADAS
% ============================================================================
\section{Fuentes y Referencias Consultadas (Criterios V7a, V7b, V7c)}

\begin{enumerate}[leftmargin=1.5em, itemsep=2pt, topsep=2pt]
    \item \textbf{Apunte Base:} Higueras Cheuquepan, N. (2026). \textit{Introducción a la integral}. MAT1188 Cálculo Intermedio, UCT.
    \item \textbf{Fuente Complementaria 1:} Stewart, J. (2012). \textit{Cálculo de una variable. Trascendentes tempranas} (7.\textsuperscript{a} ed.). Cengage Learning. ISBN 978-607-481-881-9. Secciones 5.1 (p. 360) y 5.2 (p. 371).
    \item \textbf{Fuente Complementaria 2:} Strang, G. y Herman, E. (2022). \textit{Cálculo}, vol. 1, §5.1 («Aproximación de áreas»). OpenStax. \url{https://openstax.org/books/calculo-volumen-1/pages/5-1-aproximacion-de-areas}
    \item \textbf{Referencia Visual:} Sanderson, G. [3Blue1Brown]. (2017). \textit{Integration and the fundamental theorem of calculus} [Video]. YouTube (min. 8:20--11:18). \url{https://www.youtube.com/watch?v=rfG8ce4nNh0}
    \item \textbf{Documentación Manim:} The Manim Community Developers (2026). \textit{CoordinateSystem} (v0.22.0) -- Método \texttt{get\_riemann\_rectangles}. \url{https://docs.manim.community/en/stable/reference/manim.mobject.graphing.coordinate_systems.CoordinateSystem.html}
\end{enumerate}

% ============================================================================
% 7. DECLARACIÓN DE USO DE INTELIGENCIA ARTIFICIAL
% ============================================================================
\section{Declaración de Uso de Inteligencia Artificial y Verificación (Criterios V7d, V7e)}

\textit{(Detalle completo en \href{https://github.com/sofiarodrigueznog/EID-Grupo-3/blob/main/docs/declaracion_ia.md}{\texttt{docs/declaracion\_ia.md}})}. El equipo declara el uso ético de herramientas de IA bajo verificación humana estricta:
\begin{itemize}[leftmargin=1.5em, itemsep=2pt, topsep=2pt]
    \item \textbf{Denys Sofía Rodríguez Noguera:} \textbf{Claude} (Anthropic) para arquitectura del repositorio, workflow CI (\texttt{nombre-de-rama.yml}) y pruebas con Jira. \textit{Verificación:} Pruebas manuales de PRs y revisión antes de cada merge a \texttt{main}.
    \item \textbf{Paulo Villalobos Aravena:} \textbf{GitHub Copilot} para asistencia con comandos Git y \textbf{ChatGPT} (OpenAI) para crítica de estilo y tecnicismo en \texttt{teoria.md} y \texttt{ejemplo.md}. \textit{Verificación:} Deducción manual de las sumas y contraste bibliográfico.
    \item \textbf{Antonio Ignacio Lara Fuentealba:} \textbf{ChatGPT} (OpenAI) para configuración de parámetros en \texttt{config.py}, funciones de escenas y tests en \texttt{test\_config.py}. \textit{Verificación:} Inspección visual del renderizado en Manim y ejecución de \texttt{unittest}.
    \item \textbf{Renato Misrrain Carrasco Oviedo:} \textbf{Antigravity / Gemini} (Google DeepMind) para estructura de \texttt{verify.py}, runner \texttt{render.py}, Escena 3 y fluidez del guion. \textit{Verificación:} Auditoría con fracciones exactas, compilación en 480p/1080p y validación cruzada.
\end{itemize}

\end{document}
